\begin{frame}{교란 실험 결과: 숫자가 말하는 것}
  \begin{columns}[T]
    \begin{column}{0.52\textwidth}
      \small
      \begin{tabular}{@{}lcc@{}}
        \toprule
        결과 (200회) & BC(시연 5개) & DAgger(2라운드) \\
        \midrule
        목표 도달 & 57회 (28.5\%) & \textbf{185회 (92.5\%)} \\
        위험칸 진입 & 74회 (37\%) & 13회 (6.5\%) \\
        미도달(타임아웃) & 69회 (34.5\%) & 2회 (1\%) \\
        \bottomrule
      \end{tabular}
    \end{column}
    \begin{column}{0.46\textwidth}
      \centering
      \includegraphics[width=\linewidth]{ch12_1_perturbation_study.png}
      {\footnotesize 시작 $(1,3)$에서 200회 롤아웃 — BC vs DAgger}
    \end{column}
  \end{columns}
  \vskip0.6em
  \begin{itemize}
    \item \textbf{BC의 37\%가 위험칸 진입} — 시연엔 위험칸이 하나도
      없는데, 기본값 right로 시연 밖에서 배회하다 공사 구간에 밀려든다.
    \item ``시연이 안전했으니 정책도 안전할 것'' — \textbf{시연 분포와
      배포 분포가 일치할 때만} 성립.
  \end{itemize}
\end{frame}
